\section{模型假设与符号说明}
\label{sec:assumptions}

\subsection{模型假设}
\subsubsection{全局假设}
\begin{enumerate}[label=\textbf{A\arabic*.},leftmargin=2.6em,itemsep=2pt]
  \item \textbf{损失由配置决定。}在给定模型族与训练配方下，验证集交叉熵损失 $L$ 由参数量 $N$、
    训练 token 数 $D$、数据质量 $Q$ 与领域配比 $\boldsymbol{p}$ 共同决定，
    训练过程中的超参数调整与随机波动不构成系统性的解释变量。
  \item \textbf{附件可信度分级可采信。}各附件所标注的真实观测、外推子集与半合成性质被采信；
    半合成数据只用于补充分析，其结论不表述为直接实验观测。
  \item \textbf{尺度可迁移性有限成立。}配比到损失的响应形状可跨训练尺度迁移，
    但水平项随尺度平移。因此跨尺度验证一律在按尺度去除水平偏移之后进行，或改用秩相关度量。
  \item \textbf{算力近似成立。}$C_{\text{train}}=6ND$ 作为前向与反向传播算力的近似，
    在附件 B 的真实轨迹上该口径与报告的算力值一致。
  \item \textbf{损失与得分单调负相关。}多维基准的综合得分与验证集交叉熵损失单调负相关，
    这是问题四在两种尺度之间建立映射的前提；该关系本身的拟合误差作为独立随机项处理。
  \item \textbf{时间轴以提交日期计。}问题四按模型在榜单上的提交日期归入月份，
    并把投榜模型视为同期开源前沿的代表性样本，而非开源模型的总体。
\end{enumerate}

\subsubsection{方法性假设}
除上述全局假设外，各问还依赖若干会被结论直接继承的前提，此处一并声明。

\begin{enumerate}[label=\textbf{B\arabic*.},leftmargin=2.6em,itemsep=2pt]
  \item \textbf{质量指标方向表由本队推得}（问题一）。附件未给出 22 个指标的方向标注，
    方向依据字段定义、指标间一致性与可核验的文本证据推得，并按可信度分三层；
    方向不确定的层次承重项由本队裁定，其影响以敏感性方式报告。
  \item \textbf{配比效应的形状可跨来源迁移}（问题二）。附件中没有任何一个文件同时给出
    配比与 $(N,D)$ 的联合观测，故 $b$ 的取值只能由问题一的配比模型引入。
    该假设不可在本数据上检验，且被结论直接继承。
  \item \textbf{数据量指数可跨来源共享}（问题二）。质量指数的反解依赖于该前提。
    本文对该前提单独做了检验并如实报告结果，质量指数的解释因此带有条件性。
  \item \textbf{领域配比由前两问固定}（问题三）。配比项对配比是线性的，
    在单纯形约束下最优解落于顶点，故本问把配比作为既有配置而非寻优变量。
  \item \textbf{基线质量取固定值}（问题三）。$Q_0$ 的取值未经外部确认，
    本文在其合理范围内做扫动并报告结论的稳定性。
  \item \textbf{榜单得分与算力元数据可经名称连接}（问题四）。
    两个附件中的模型经名称规范化后视为同一模型，该连接是承重的；
    名称无法匹配的记录作缺失处理，不参与分解。
\end{enumerate}

\subsection{符号说明}
\subsubsection{全局符号}
表~\ref{tab:symbols-global}列出全文通用符号。凡涉及同一符号在不同子问题中含义不同的情形，
本文在首次出现处以下标或修饰符区分，并在表中注明。

\begin{longtable}{@{}p{0.13\linewidth}p{0.24\linewidth}p{0.20\linewidth}p{0.27\linewidth}@{}}
  \caption{全文通用符号与单位}
  \label{tab:symbols-global}\\
  \toprule
  符号 & 含义 & 单位 & 说明 \\
  \midrule
  \endfirsthead
  \multicolumn{4}{@{}l}{\small 续表~\ref{tab:symbols-global}}\\
  \toprule
  符号 & 含义 & 单位 & 说明 \\
  \midrule
  \endhead
  \bottomrule
  \endfoot
  $N$ & 模型参数量 & 十亿（B） & 全文以十亿为单位；代入 $C=6ND$ 时须乘以 $10^9$ \\
  $D$ & 训练 token 数 & 十亿（B） & 同上 \\
  $C$ & 算力预算 & FLOPs & 问题三沿 $10^{18}$--$10^{25}$ 扫描，以 $10^{19}$、$10^{22}$、$10^{24}$ 三档锚定 \\
  $L$ & 验证集交叉熵损失 & nats/token & 全文的损失尺度 \\
  $Q$ & 数据质量评分 & 无量纲 & 取值于 $(0,1]$；问题一给出，问题二、三沿用 \\
  $Q_0$ & 基线质量 & 无量纲 & 问题三中不投入质量提升时的质量水平 \\
  $\boldsymbol{p}$ & 领域配比向量 & 无量纲 & $\boldsymbol{p}=(p_1,\dots,p_{17})$，$p_i\ge 0$ 且 $\sum_i p_i=1$ \\
  $E$ & 不可约损失 & nats/token & 经典标度律的截距项 \\
  $A,\ \alpha$ & 参数量项的系数与指数 & — & 由真实轨迹拟合得到 \\
  $B,\ \beta$ & 数据量项的系数与指数 & — & $\beta$ 亦记作数据量弹性 \\
  $\gamma$ & 质量指数 & 无量纲 & $D_{\text{eff}}=D\,Q^{\gamma}$ 中的指数 \\
  $D_{\text{eff}}$ & 有效 token 数 & 十亿（B） & $D_{\text{eff}}=D\,Q^{\gamma}$，把质量折算为等效数据量 \\
  $\boldsymbol{b}$ & 配比项的系数向量 & 无量纲 & $\boldsymbol{b}=\kappa(\boldsymbol{\theta}-\bar\theta)$，$\kappa<0$，在附件 B 上不可辨识 \\
  $\boldsymbol{z}_{\boldsymbol{p}}$ & 配比的中心化偏离 & 无量纲 & 均匀配比对应 $\boldsymbol{z}_{\boldsymbol{p}}=\boldsymbol{0}$ \\
  $\kappa$ & 配比项的标度常数 & — & 在附件 B 上不可辨识，故只报比值 $b_i/b_j$ \\
    $g(Q)$ & 单位数据质量的提升成本函数 & — & 问题三给出三类候选形式 \\
  $g_{\mathrm{c}},\ g_{\mathrm{r}}$ & 成本函数的幅值系数与速率参数 & — & 仅在问题三讨论第四类 $g$ 时使用，避免与质量指数 $\gamma$ 同名 \\
  $L_{\text{ctx}}$ & 上下文窗口长度 & token & 外生给定，取值依据附件 C7 \\
  $\eta$ & 注意力开销系数 & — & $\eta=2\times10^{-4}$ \\
  $k$ & 单位数据的训练与注意力开销系数 & — & $k=6+\eta L_{\text{ctx}}$ \\
  $k_{\mathrm{H}}$ & Huber $M$ 估计的稳健化参数 & — & 取 $1.345$，仅用于问题一的域级聚合 \\
  $B_{\mathrm{boot}}$ & 自助重抽样次数 & — & 仅用于问题一的区间估计，与数据量项系数 $B$ 无关 \\
  $\tau$（问题一，无下标） & 域序的 Kendall 秩相关系数 & 无量纲 & 仅用于问题一的域序稳定性；问题四的 $\tau_t$ 是月度效应，二者不同 \\
  $Q^{*},N^{*},D^{*}$ & 最优配置 & 无量纲 / B / B & 问题三在给定预算下的解 \\
  $S$ & 模型综合能力得分 & 分 & 仅用于问题四，为六维基准的等权平均 \\
  $S_{q90}$ & 能力前沿 & 分 & 每月综合能力的 $0.90$ 分位 \\
  $N_{q90}$ & 前沿模型规模 & 十亿（B） & 每月参数量的 $0.90$ 分位 \\
  $f(N)$ & 规模通道函数 & 分 & 问题四分解中的规模项 \\
  $\tau_t$ & 第 $t$ 月的技术通道效应 & 分 & 吸收规模之外的全部月度共同变动 \\
  $V,\ P$ & 价值维度、污染维度得分 & 无量纲 & 问题一门控合成的两个分量 \\
  $\theta_d$ & 第 $d$ 个配方域的损失水平系数 & nats/token & 问题一配比模型的系数；问题二的 $\boldsymbol{b}$ 由其中心化得到 \\
\end{longtable}

\subsubsection{数据与附件编号}
附件编号沿用《数据说明》：A1--A3 为质量信号的全量记录，A4--A11 为三尺度配方与损失，
A12--A15 为高尺度外推表，A16 为跨体系映射参考，A17、A18 为辅助摘要与抽样文本统计；
B1--B12 为标度律轨迹、收敛点、文献值与半合成网格；C1--C10 为榜单、模型元数据、
桥接数据与逐任务评测明细。四问的附件映射见第~\ref{sec:problem}~节图~\ref{fig:data-pipeline}。

\subsubsection{指标与领域的简称}
质量侧的 22 个指标与配比侧的 17 个领域均以附件原名引用。为便于行文，
表~\ref{tab:symbols-abbrev}给出正文中反复出现的少数简称。

\begin{table}[H]
  \centering
  \caption{正文使用的简称对照}
  \label{tab:symbols-abbrev}
  \begin{tabular}{@{}lp{0.76\linewidth}@{}}
    \toprule
    简称 & 全称或含义 \\
    \midrule
    门控合成 & 质量评分 $Q=V\cdot(1-P)$，$V$ 为价值维度均值，$P$ 为污染维度污染度均值 \\
    可分辨域对 & 两个领域的质量区间（自助 95\%）不重叠，故其高低可在统计上区分 \\
    冲突率 & 由所采用的冲突判据判定为冲突的样本占全体的比例 \\
    去偏 $R^2$ & 在各尺度上先扣除水平偏移后再计算的决定系数 \\
    去源偏 RMSE & 先扣除各来源的截距偏移后再计算的均方根误差 \\
    端点法 & 仅用窗口首月与末月的读数做差分以定义通道贡献 \\
    趋势法 & 用窗口内全部十个月的斜率定义通道贡献 \\
    活跃约束集 & 在最优解处取等号的约束集合 \\
    结构性转移 & 最优策略随预算发生质变；其数学定义见第~\ref{sec:problem3}~节 \\
    \bottomrule
  \end{tabular}
\end{table}
